\section{Systematic Uncertainties and Result}
\label{sec:syst}
\emph{In short: thirteen variants in four groups give a systematic uncertainty of $0.0070$ (Bayesian) and $0.0072$ (frequentist), dominated by the particle-identification requirement; the final result is $\Araw=-0.0443\pm0.0085\,(\text{stat})\pm0.0070\,(\text{syst})$.}
The frequentist and Bayesian fits share one systematic study: the same thirteen variants of the nominal choices, in four groups, and the same rule to combine them. Each variant is fitted by both methods on exactly the same candidates. The variants are: \emph{window}, lower edge at $5050$ or $5150\,\MeV$ and upper edge at $5450$ or $5700\,\MeV$; \emph{background}, a single exponential per charge, a double exponential with the shape shared between charges, or a charge-symmetric background ($A_b=0$); \emph{signal}, a Crystal Ball shape with free tail parameters (Eq.~\eqref{eq:cb}) or $\mu$ and $\sigma$ separate for the two charges; \emph{preselection}, one cut at a time loosened or tightened, $\texttt{ProbK}>0.4$ or $>0.6$ and $\texttt{ProbPi}<0.4$ or $<0.6$, on all three tracks.

The shift $\Delta=\Araw^{\mathrm{var}}-\Araw^{\mathrm{nom}}$ of a variant is not necessarily a systematic effect: when a variant changes the sample (window, preselection) the two samples are nested and part of $\Delta$ is a statistical fluctuation of expected size $\sigma_B=\sqrt{|\sigma_{\mathrm{var}}^2-\sigma_{\mathrm{nom}}^2|}$~\cite{barlow}, and such a variant contributes only if $|\Delta|>\sigma_B$. A variant that the data reject ($\Delta\chi^2>25$ on the same data) is excluded: this removes the single exponential ($\Delta\chi^2=134$, Appendix~\ref{sec:bayes_gof}). The largest remaining $|\Delta|$ of each group is its uncertainty, and the groups are added in quadrature.

\begin{table*}[t]
\centering
\resizebox{\ifdim\width>\linewidth\linewidth\else\width\fi}{!}{\begin{tabular}{llrrrrr}
\toprule
Group & Variant & $n$ & $\Delta$ Bayes & $\sigma_B$ Bayes & $\Delta$ freq. & $\sigma_B$ freq. \\
\midrule
Window & $[5050,5600]$ & $21164$ & $-0.0003$ & $0.0009$ & $-0.0004$ & $0.0005$ \\
 & $[5150,5600]$ & $18566$ & $-0.0001$ & $0.0007$ & $-0.0005$ & $0.0021$ \\
 & $[5100,5450]$ & $19004$ & $+0.0002$ & $0.0009$ & $-0.0001$ & $0.0003$ \\
 & $[5100,5700]$ & $20191$ & $\mathbf{-0.0014}$ & $0.0010$ & $\mathbf{-0.0014}$ & $0.0001$ \\
Background & single exponential (rejected) & $19731$ & $-0.0037$ & -- & $-0.0039$ & -- \\
 & shape shared between charges & $19731$ & $\mathbf{-0.0037}$ & -- & $\mathbf{-0.0040}$ & -- \\
 & symmetric, $A_b=0$ & $19731$ & $-0.0028$ & -- & $-0.0031$ & -- \\
Signal & Crystal Ball & $19731$ & $\mathbf{-0.0002}$ & -- & $\mathbf{-0.0006}$ & -- \\
 & $\mu,\sigma$ per charge & $19731$ & $-0.0000$ & -- & $-0.0001$ & -- \\
Preselection & $\texttt{ProbK}>0.4$ & $26553$ & $-0.0032$ & $0.0038$ & $-0.0032$ & $0.0037$ \\
 & $\texttt{ProbK}>0.6$ & $13981$ & $\mathbf{+0.0058}$ & $0.0048$ & $\mathbf{+0.0058}$ & $0.0050$ \\
 & $\texttt{ProbPi}<0.4$ & $17984$ & $+0.0022$ & $0.0022$ & $+0.0020$ & $0.0023$ \\
 & $\texttt{ProbPi}<0.6$ & $20953$ & $-0.0004$ & $0.0020$ & $-0.0006$ & $0.0017$ \\
\midrule
\multicolumn{3}{l}{Group uncertainties: window / background / signal / preselection} & \multicolumn{2}{c}{$0.0014/0.0037/0.0002/0.0058$} & \multicolumn{2}{c}{$0.0014/0.0040/0.0006/0.0058$} \\
\multicolumn{3}{l}{Total (quadrature)} & \multicolumn{2}{c}{$0.0070$} & \multicolumn{2}{c}{$0.0072$} \\
\bottomrule
\end{tabular}}
\caption{Systematic study common to the Bayesian and frequentist fits: shift of $\Araw$ for each variant and, for the variants that change the sample, the expected statistical spread $\sigma_B$ of the shift. The largest contributing shift of each group is in bold. The single exponential is rejected by the data and excluded. For the window $[5150,5600]$ the frequentist uncertainty is scaled from the number of events, since the Hesse matrix of that fit is not accurate.}
\label{tab:syst}
\end{table*}

The results are in Table~\ref{tab:syst}, and all the Bayesian fits in Table~\ref{tab:bayes_all} of the Appendix. The two methods give the same shifts to within $\sim$$0.0005$ and the same group uncertainties, with totals of $0.0070$ (Bayesian) and $0.0072$ (frequentist):
\begin{gather*}
\text{Bayesian:}\\ \Araw=-0.0443\pm0.0085\,(\text{stat})\pm0.0070\,(\text{syst}),\\[2pt]
\text{frequentist:}\\ \Araw=-0.0441\pm0.0085\,(\text{stat})\pm0.0072\,(\text{syst}).
\end{gather*}
The preselection dominates, through the single variant $\texttt{ProbK}>0.6$ ($+0.0058$, $1.2\,\sigma_B$), a statistical fluctuation of a nested sample: a scan of each cut in steps of $0.05$ shows no trend on the tight side (all within $1.3\,\sigma_B$), while loosening the kaon requirement moves $\Araw$ monotonically, as expected if misidentified decays not in the mass model enter the sample, which supports the nominal requirement. The background group is second, driven by the charge dependence of the background, and the window third.

\subsection{Stability Checks of the Frequentist Fit}
\label{sec:freq_syst}
\label{sec:syst_scans}
\emph{In short: the fit does not depend on the fit window or on the signal region, and the only dependence that matters is the one on the particle-identification requirements.}
While the systematic uncertainties study is shared with the Bayesian approach -- and discussed in detail in Sec.~\ref{sec:syst} -- three specific stability checks were performed within the frequentist one, to evaluate the robustness of the fitted raw asymmetry $\Araw$. Each test consists of a two-dimensional grid scan in which two selection or fit parameters are varied simultaneously, performing a simultaneous fit with the nominal double-exponential background model at each point and observing the stability of $\Araw$.

\begin{enumerate}[leftmargin=*]
    \item \textbf{Fit Window Scan:} The lower and upper boundaries of the invariant mass fit window, nominally set to $[5100, 5600]\,\MeV$, were scanned over the ranges $[5000, 5200]\,\MeV$ and $[5350, 6100]\,\MeV$, respectively, with a 50 \MeV step. This test evaluates the sensitivity of $\Araw$ to the exponential background parametrization, and is useful to choose the nominal window: too small or too asymmetric fit window produces unstable fit, due to worse background estimation.

    \item \textbf{Signal Region Scan:} Changing the extremes of the signal region not significant fluctation in $\Araw$ are observed, suggesting that the sideband preliminar fit could be avoided.

    \item \textbf{Particle Identification Cut Scan:} The preselection cuts on the final-state kaon PID likelihoods (\texttt{ProbK} and \texttt{ProbPi}) were varied simultaneously towards a better kaon purity (larger than $0.5$ for kaons, smaller than $0.5$ for pions). This influences significantly the $\Araw$ estimation, but always within the $68\%$ CI, confirming the cut selection as reasonable.
\end{enumerate}

Across all scan configurations, the maximum observed variation in $\Araw$ remains below the statistical uncertainty ($\Delta \Araw < \sigma_{\mathrm{stat}}$), demonstrating the stability of the nominal fit result.

\paragraph*{Numbers.} Three results from the same notebook, which do not enter the systematic uncertainty. Over the grid of window edges the maximum shift is $0.0048$ (mean statistical uncertainty $0.0085$); the signal region, used only to start the sideband fit, moves $\Araw$ by at most $6\times10^{-5}$; on the grid of the two PID cuts the maximum shift is $0.0126$, larger than the nominal $0.0085$ but within the statistical uncertainty of the much smaller samples ($0.0112$ on average), and no shift exceeds $1.3\,\sigma_B$ in the one-cut-at-a-time scan on the tight side. The preselection therefore remains the dominant contribution.